% 5.振动波动.tex 子文件，book文档类，主文档通过\input{5.振动波动.tex}导入
\part{第三篇 振动与波动}
\setcounter{chapter}{10}
\chapter{第十一章 振动学基础}

\section{简谐运动的描述}
    \subsection{弹簧振子}
    \subsection{简谐运动的表达式}
    \subsection{简谐运动的速度和加速度}
    \subsection{振动的相位}
    \subsection{旋转矢量表示法}

\section{简谐振动的动力学特征}
    \subsection{简谐振动的动力学定义}
    \subsection{简谐振动的实例}
    \subsection{简谐振动的能量}

\section{简谐振动的合成}
    \subsection{两个同方向同频率简谐振动的合成}
    \subsection{两个同方向不同频率简谐振动的合成}
    \subsection{相互垂直的简谐振动的合成}
    质点的轨迹方程为
    %\[
    %\frac{x^2}{A_1^2}+\frac{y^2}{A_2^2}-\frac{2xy}{A_1A_2}\cos\Delta\varphi=\sin^2\Delta\varphi
    %\]

\section{阻尼振动}

\section{受迫振动和共振}

\section{电磁振荡}

\chapter{第十二章 波动学基础}

\section{机械波的产生和传播}
    \subsection{机械波的形成}
    \subsection{描述波动的物理量}

\section{平面简谐波的波函数}
    \subsection{波函数的建立}
    写成下列形式：
    \[
    y(x,t)=A\cos\left[2\pi\left(\frac{t}{T}-\frac{x}{\lambda}\right)\right]
    \]
    我们，其定义是
    %\[
    %\]
    它可以理解为，波函数式又可写成
    %\[
    %\]
    如果

    以上这些，且，至于，即已知的振动表达式为
    %\[
    %y()
    %\]
    则相应的在$x$轴上传播的简谐波的波函数为
    \[
    y(x,t)=A\cos{\left[\omega\left(t+\frac{x-x_0}{u}\right)+\varphi\right]}
    \]
    \[
    y = A \cos\left[\omega t + \varphi_0 \mp \frac{2\pi}{\lambda} (x - d)\right]
    \]
    \subsection{波函数的物理意义}

\section{波动方程与波速}
    \subsection{物体的弹性形变}
        \subsubsection{线变}
        \subsubsection{切变}
    \subsection{波动方程}
    \subsection{决定波速的因素}

\section{波的能量}
    \subsection{平面简谐纵波传播时介质元的能量}
    仍，当
    %\[
    %\]
    在介质，速度为
    %\[
    %\]
    则介质元的振动动能为
    %\[
    %dE_k=\frac{1}{2}\rho dV\omega^2 A^2\sin^2\left[\omega\left(t-\frac{x}{u}\right)\right]
    %\]

    介质，故弹性势能可写成
    %\[
    %\]
    \subsection{波的能量密度和能流密度}
    介质，可得
    %\[
    %\]
    式说明，因为，所以能量密度的平均值为
    \[
    \overline{w}=\frac{1}{T}\int^T_0 w\mathrm{d}t=\frac{1}{2}\rho A^2 \omega^2
    \]
    式表示波的

    因为，数值$I$为
    \[
    I=\frac{\omega uS}{S}=u=\frac{1}{2}\rho A^2\omega^2 u
    \]
    能流密度

    由于
    \subsection{波的吸收}
    即
    %\[
    %\]
    式，可得
    %\[
    %\]
    式，为
    %\[
    %\]
    \subsection{球面波}

\section{惠更斯原理}
    \subsection{惠更斯原理}
    \subsection{惠更斯原理的应用}

\section{波的叠加原理和波的干涉}
    \subsection{波的叠加原理}
    \subsection{波的干涉}

\section{驻波}
    \subsection{驻波的形成}
    \subsection{驻波的波函数}
    下面，则它们的波函数分别为
    %\[
    %y_1=
    %\]
    %\[
    %y_2=
    %\]
    两波叠加后，合位移为
    \[
    y=y_1+y_2=A\cos\left[2\pi\left(\nu t-\frac{x}{\lambda}\right)\right]+A\cos\left[2\pi\left(\nu t+\frac{x}{\lambda}\right)\right]
    \]
    利用得
    \[
    y=\left(2A\cos{\frac{2\pi}{\lambda}x}\right)\cos{2\pi \nu t}
    \]
    这就是

    从

    \subsection{半波损失}
    波，当波从波疏介质垂直入射波密有半波损失。
    \subsection{弦线振动的简正模式}

\section{多普勒效应}
    现以，$u$表示声波在介质中的传播速度

    设波源频率为，下面分三种情况讨论：

    (1) 波源$S$静止在介质中，观测者相对于介质以速度$v_0$向着波源移动，观测者向着波源运动时，感到波是以速度为，即观测者测得的频率为
    %\[
    %\]

    (2) 设观测者静止在介质中，看出
    %\[
    %\]
    所以观测者每秒内接收到的完整波的个数，即频率$\nu'$为
    \[
    \nu'=\frac{u}{\lambda'}=\frac{u}{(u-u_s)T}=\nu\cdot\frac{u}{u-u_s}
    \]
    当

    (3) 设波源，为
    \[
    \nu'=\nu\left(\frac{u+u_0}{u-u_s}\right)
    \]

\section{声波}

\section{电磁波}

\endinput